% Modern editorial, markup and figure/layout contributions: CC-BY-SA-4.0.
% Historical text and illustrations: Leonhard Euler (1744), public domain.
% Component scope and upstream attribution: see RIGHTS.md.
% Image-collated Latin transcription, source PDF230–247.
\subsection*{Caput VI.}
\textit{Methodus, inter omnes curvas pluribus proprietatibus communibus gaudentes, eam determinandi quæ maximi minimive proprietate sit prædita.}
\subsection*{Propositio I. Theorema.}
\originalsection{1}\sourcepage{227}\textit{Curva, quæ inter omnes omnino curvas habet expressionem $\alpha A+\beta B$ maximam vel minimam, eadem simul ita erit comparata, ut inter omnes eadem proprietate $A$ præditas contineat valorem formulæ $B$ maximum vel minimum.}
\subsection*{Demonstratio.}
Ponamus inventam esse curvam, in qua inter omnes alias eidem abscissæ respondentes valor expressionis $\alpha A+\beta B$ sit maximus; quod enim de maximo demonstrabitur, idem mutatis mutandis de minimo valebit. Denotant autem litteræ $A$ \& $B$ hic nobis ejusmodi formulæ vel expressiones indeterminatæ, in quas Quæstio de maximis \& minimis cadere queat; tum vero \sourcepage{228}$\alpha$ \& $\beta$ sunt quantitates constantes quæcunque. Designemus jam istam curvam, in qua sit $\alpha A+\beta B$ maximum, littera $Q$, quo eam facilius sine molesta verborum descriptione indicare queamus; Nunc concipiatur alia quæcunque curva $R$ eidem abscissæ respondens, quæ recipiat formulæ $A$ eundem valorem, quem tenet curva $Q$: in hac igitur curva $R$ expressio $\alpha A+\beta B$ minorem occupabit valorem, quam in curva $Q$; eo quod in curva $Q$ expressio $\alpha A+\beta B$ omnium maximum valorem sortitur. Quare cum in curvis $Q$ \& $R$ expressio $A$ eundem obtineat valorem, atque in $Q$ expressio $\alpha A+\beta B$ major sit quam in curva $R$; sequitur in curva $Q$ valorem expressionis $B$ majorem esse debere quam in curva $R$. Cum igitur $R$ curvam quamcunque denotet, quæ cum $Q$ communem valorem formulæ $A$ recipiat; manifestum est inter omnes has curvas $R$ curvam $Q$ esse illam, in qua formula $B$ maximum habeat valorem. Ex quibus conficitur, eam curvam, quæ inter omnes omnino curvas habeat expressionis $\alpha A+\beta B$ valorem maximum vel minimum, eandem curvam simul ita esse comparatam, ut inter omnes alias curvas secum eadem communi proprietate $A$ gaudentes possideat maximum minimumve valorem expressionis $B$. Quanquam enim Demonstratio tantum ad maximum est adornata, tamen eadem, translatis verbis, ad minimum accommodabitur. Q. E. D.
\subsection*{Corollarium I.}
\originalsection{2}Vicissim itaque intelligitur, si curva debeat investigari, quæ inter omnes alias eadem communi proprietate $A$ præditas expressionem $B$ sit habitura maximam vel minimam; tum quæsito satisfieri, si absolute inter omnes curvas ea definiatur, in qua sit $\alpha A+\beta B$ maximum vel minimum.
\subsection*{Corollarium II.}
\originalsection{3}In solutionem igitur hujusmodi Problematum binæ novæ ingredientur constantes arbitrariæ $\alpha$ \& $\beta$, quæ in ipsis expressionibus \sourcepage{229}$A$ \& $B$ non inerant: hæ autem unius duntaxat constantis vicem sustinebunt; quia earum ratio tantum in computum venit.
\subsection*{Corollarium III.}
\originalsection{4}Quod si ergo, inter omnes curvas eadem communi proprietate $A$ gaudentes, eam definiri oporteat in qua sit $B$ maximum minimumve; tum utriusque expressionis $A$ \& $B$ capiantur valores differentiales, qui, per constantes arbitrarias seorsim multiplicati \& conjunctim nihilo æquales positi, dabunt æquationem pro curva quæsita.
\subsection*{Corollarium IV.}
\originalsection{5}Simul etiam perspicuum est perinde esse, sive inter omnes curvas eadem communi proprietate $A$ gaudentes, ea quæratur in qua sit $B$ maximum vel minimum; sive vicissim inter omnes curvas eadem communi proprietate $B$ gaudentes, ea quæratur in qua sit $A$ maximum vel minimum.
\subsection*{Scholion.}
\originalsection{6}Quæ, cum in hac Propositione, tum in annexis Corollariis, tradidimus, ex Capite præcedente jam sunt planissima: quippe quibus continetur inversa Methodus resolvendi Problemata, in quibus, inter omnes curvas eadem communi proprietate gaudentes, ea quæritur quæ prædita sit maximi minimive alicujus indole. Neque vero idcirco idem argumentum nos tantum repetivisse censendum est; nam eandem veritatem, quam ante modo satis prolixo elicueramus, hic admodum succincte \& breviter dedimus demonstratam. Quocirca eo fortius altera demonstrandi Methodus per alteram confirmabitur, ob summum utriusque consensum: atque si cui prior Methodus non satis perfecta, propter tantam infinite parvorum compagem, nimis lubrica \& incerta videatur, ei Demonstratio hic data omnem
\sourcepage{230} scrupulum adimet. Deinde, si quis de præsentis Propositionis conversione in Coroll. 1 facta etiamnum dubitet, ei prior Methodus plenissime satisfaciet. Interim ratio conversionis ex se satis tuto inferri potest. Cum enim curva $Q$, quæ inter omnes omnino curvas habeat $\alpha A+\beta B$ maximum vel minimum, ita sit comparata, ut inter omnes curvas eadem communi proprietate $A$ gaudentes, habeat $B$ maximum vel minimum, quicquid loco $\alpha$ \& $\beta$ accipiatur; necesse est ut conversio æque pateat, siquidem coefficientibus $\alpha$ \& $\beta$ summa extensio tribuatur. Hocque adeo commemorare, hujusque ratiocinii validitatem declarare visum est, ut in sequentibus, ubi eodem utemur, nullum dubium relinquatur. Hanc enim Propositionem, etsi proprie ad Caput præcedens pertinet, huc transtulimus, quo eadem Methodo proprium hujus Capitis argumentum facilius pertractare possimus; quippe quod, si altera Methodo expediri deberet, prolixissimos requireret calculos, maximasque differentialium omnium ordinum tricas. Interim tamen, quantum fieri potest, dilucide ostendemus omnia, quæ hic trademus, per Methodum superiorem confirmari atque etiam elici posse.
\subsection*{Propositio II. Theorema.}
\originalsection{7}\textit{Quæ curva, inter omnes omnino curvas eidem abscissæ respondentes, habet valorem expressionis $\alpha A+\beta B+\gamma C$ maximum vel minimum, eadem curva simul ita erit comparata, ut inter omnes curvas, quæ tam expressionem $A$ quam expressionem $B$ communem habent, possideat valorem expressionis $C$ maximum vel minimum.}
\subsection*{Demonstratio.}
Denotant hic nobis litteræ $A$, $B$ \& $C$ formulas integrales vel expressiones indefinitas ejusmodi, quæ maximi minimive sint capaces; at litteræ $\alpha,\beta,\gamma$ designant quantitates constantes arbitrarias. Sit nunc $Q$ curva, quæ inter omnes omnino curvas \sourcepage{231}habeat valorem $\alpha A+\beta B+\gamma C$ maximum vel minimum; atque concipiatur alia quæcunque curva $R$, in qua, cum expressio $A$, tum $B$, eundem obtineat valorem quem obtinet in curva $Q$; quo posito expressio composita $\alpha A+\beta B$ eundem habebit valorem in utraque curva $Q$ \& $R$. Hanc ob rem, expressio tota $\alpha A+\beta B+\gamma C$ in curva $R$ minorem sortietur valorem quam in curva $Q$, siquidem $\alpha A+\beta B+\gamma C$ in curva $Q$ est maximum: contra expressionis $\alpha A+\beta B+\gamma C$ valor in curva $R$ major erit quam in curva $Q$, si $\alpha A+\beta B+\gamma C$ in curva $Q$ fuerit minimum. Cum igitur expressionis portio $\alpha A+\beta B$ utrique curvæ $Q$ \& $R$ sit communis, reliqua portio $\gamma C$, atque adeo expressio $C$, in casu maximi major erit in $Q$ quam in $R$, in casu minimi autem expressio $C$ in curva $Q$ minor erit quam in curva $R$. Ex quibus sequitur, si curva $Q$ inter omnes omnino curvas, habuerit valorem expressionis $\alpha A+\beta B+\gamma C$ maximum vel minimum, tum simul hanc curvam $Q$ ea indole esse præditam, ut inter omnes curvas $R$ quæ eodem valore cum expressionis $A$ tum expressionis $B$ gaudeant, contineat valorem expressionis $C$ maximum vel minimum. Q. E. D.
\subsection*{Corollarium I.}
\originalsection{8}Quoniam expressiones $A$, $B$ \& $C$ pro lubitu inter se commutari possunt; curva in qua est $\alpha A+\beta B+\gamma C$ maximum vel minimum, ea simul vel, inter omnes curvas iisdem proprietatibus $A$ \& $B$ communibus gaudentes, habebit $C$ maximum vel minimum; vel habebit $B$ maximum minimumve, inter omnes curvas quæ proprietatibus $A$ \& $C$ communibus gaudebunt; vel denique habebit $A$ maximum minimumve, inter omnes curvas in quas ambæ proprietates $B$ \& $C$ æque competent.
\subsection*{Corollarium II.}
\originalsection{9}Quæ igitur curva, inter omnes iisdem binis proprietatibus $A$ \& $B$ communibus gaudentes, habet $C$ maximum minimumve; \sourcepage{232}eadem habebit inter omnes curvas binis proprietatibus vel $A$ \& $C$, vel $B$ \& $C$, æque præditas, vel $B$, vel $A$ maximum minimumve.
\subsection*{Corollarium III.}
\originalsection{10}Si igitur curva quæri debeat quæ, inter omnes alias binis proprietatibus $A$ \& $B$ æqualiter præditas, habeat expressionem $C$ maximam vel minimam; tum quæsito satisfiet, si curva quæratur, quæ, absolute inter omnes curvas, habeat expressionem $\alpha A+\beta B+\gamma C$ maximam vel minimam.
\subsection*{Corollarium IV.}
\originalsection{11}Quoniam $\alpha,\beta,\gamma$ sunt quantitates constantes arbitrariæ; in solutionem hujusmodi Problematum tres novæ quantitates arbitrariæ ingredientur, quæ in formulis propositis $A$, $B$, \& $C$ non inerant: æquivalebunt autem hæ tres constantes $\alpha,\beta$, \& $\gamma$ tantum duabus.
\subsection*{Corollarium V.}
\originalsection{12}Hæ vero constantes adeo jam in æquatione pro curva primum inventa inerant; præter eas vero, per integrationes novæ ingredientur constantes tot, quot integrationibus opus est, antequam ad æquationem finitam perveniatur.
\subsection*{Corollarium VI.}
\originalsection{13}Simili modo, quo hanc Propositionem \& præcedentem demonstravimus, ostendetur, curvam, quæ absolute, inter omnes curvas, habeat expressionem $\alpha A+\beta B+\gamma C+\delta D$ maximam vel minimam, eandem, inter omnes curvas tres expressiones $A$, $B$ \& $C$ communes habentes, habituram esse quartam $D$ maximam vel minimam.
\subsection*{Scholion.}
\originalsection{14}\sourcepage{233}Ex hac Propositione jam satis percipitur Methodus resolvendi ejusmodi Problemata ad Methodum relativam pertinentia, in quibus quæritur curva quæ, inter omnes eidem abscissæ respondentes \& duabus pluribusve proprietatibus communibus æque gaudentes, habeat valorem cujuspiam expressionis maximum minimumve. Quæstio scilicet perpetuo revocabitur ad Methodum absolutam; ita ut, inter omnes omnino curvas, quærenda sit curva quæ expressionem quampiam habeat maximam vel minimam. Hacque reductione id commodi nanciscimur, ut omnia hujusmodi Problemata, ope valorum differentialium quos jam supra investigare docuimus, resolvere queamus. Ipse autem resolvendi modus eo redibit, ut omnes proprietates communes, una cum maximi minimive expressione, seorsim explicentur; singulæ per constantes arbitrarias multiplicentur; \& producta in unam summam colligantur: quo facto, absolute inter omnes curvas, eam quæri oportebit, in qua ista summa sit maxima vel minima. Hoc vero ipsum perficietur, dum summæ illius valor differentialis investigabitur, nihiloque æqualis ponetur. Quocirca universa operatio absolvetur, si, cum singularum expressionum proprietates communes continentium, tum maximi minimive expressionis valores differentiales, secundum regulas supra datas, capiantur; singuli seorsim in constantes arbitrarias ducantur; omniumque horum productorum aggregatum nihilo æquale ponatur: ex quo orietur æquatio pro curva quæsita. Sufficere itaque posset hoc unicum præceptum ad Quæstiones hujus generis solvendas. Verum, antequam hujus usum exponamus, hanc ipsam Methodum via ante adhibita confirmari conveniet.
\subsection*{Propositio III. Problema.}
\originalsection{15}\textit{Inter omnes curvas ad eandem abscissam relatas, quæ binis proprietatibus communibus $A$ \& $B$ æqualiter sint præditæ; definire eam, in qua sit valor expressionis $C$ maximus vel minimus.}
\subsection*{Solutio.}
\sourcepage{234}Ex præcedentibus jam intelligitur hoc Problema solvi, si, inter omnes curvas, absolute quæratur ea in qua sit $\alpha A+\beta B+\gamma C$ maximum vel minimum. Ad hoc autem nosse oportet valores differentiales expressionum $A$, $B$, \& $C$. Sit igitur valor differentialis expressionis $A=nv.dx.P$; expressionis $B=nv.dx.Q$; expressionis $C=nv.dx.R$; ex quibus æquatio pro curva desiderata erit $\alpha P+\beta Q+\gamma R=0$.

Verum, quo hujus Solutionis veritas magis eluceat, idem hoc Problema eadem Methodo, qua supra in Capite præced. usi sumus, aggrediamur. (Fig.~15.) Primum autem intelligitur, ad hoc Problema resolvendum, ternas applicatas particulis infinite parvis augeri debere, ut tribus conditionibus præscriptis satisfieri possit. Primo enim tres has particulas adjunctas, quibus ipsa curva satisfaciens $az$ in novam a se quam minime discrepantem transmutatur, ita comparatas esse oportet, ut expressio $A$, quæ unam proprietatem communem continet, in utramque curvam æqualiter competat: Deinde etiam altera proprietas communis $B$ in utraque curva eundem valorem obtinere debebit. Tertio ex maximi minimive natura expressio quoque $C$ eundem valorem in ipsa curva \& eadem mutata nancisci debet; quibus tribus conditionibus, per pauciores quam tres particulas tribus applicatis adjunctas, satisfieri non potest. Quare præter binas applicatas $Nn$ \& $Oo$, quæ in figura particulis $nv$ \& $o\omega$ sunt auctæ, concipiatur sequenti applicatæ $Pp$ particula $p\pi$ adjici. Ac quæratur primum incrementum, quod expressio $A$ ex his tribus particulis assequitur quod erit $=nv.Pdx+o\omega.P'dx+p\pi.P''dx$. Namque ex particula $nv$ nascitur incrementum $nv.Pdx$, congruens cum ipso valore differentiali, quem expressio $A$ ex sola particula $nv$ adipiscitur. Ex sequenti vero particula $o\omega$ oritur incrementum $o\omega.P'dx$, scilicet idem quod ante, suo differentiali auctum: quia enim $o\omega$ sequenti applicatæ adjungitur, omnes quantitates $o\omega$ afficientes erunt sequentes earum, quibus particula $nv$ afficitur: atque simili ratione ex particula $p\pi$ nascetur \sourcepage{235}incrementum $p\pi.P''dx$; quæ omnia, si cui libuerit calculum eo modo quo in Cap. præc. Propos. 3 §. 22 usi sumus persequi, satis fient manifesta ac perspicua. Eodem igitur porro modo expressio $B$, cujus valorem differentialem ex unica particula $nv$ oriundum posuimus $=nv.Qdx$, ex tribus particulis $nv$, $o\omega$ \& $p\pi$ incrementum accipiet $=nv.Qdx+o\omega.Q'dx+p\pi.Q''dx$. Tertio expressio $C$ ex his tribus particulis augmentum capiet hoc $nv.Rdx+o\omega.R'dx+p\pi.R''dx$. Singula jam hæc tria incrementa seorsim nihilo æqualia poni oportet, ut omnibus conditionibus præscriptis satisfiat; unde tres sequentes æquationes orientur, facta divisione per $dx$,
\[\begin{aligned}
0&=nv.P+o\omega.P'+p\pi.P''\\
0&=nv.Q+o\omega.Q'+p\pi.Q''\\
0&=nv.R+o\omega.R'+p\pi.R''.
\end{aligned}\]
Quod si nunc particulæ $nv,o\omega,p\pi$, ad solutionem peragendam tantum in subsidium vocatæ eliminentur; orietur æquatio inter quantitates curvæ proprias, quibus proin natura curvæ exprimetur. Ad has autem particulas eliminandas, singulas æquationes per novas incognitas $\alpha,\beta,\gamma$ seorsim multiplicemus, ut habeatur
\[\begin{aligned}
0&=nv.\alpha P+o\omega.\alpha P'+p\pi.\alpha P''\\
0&=nv.\beta Q+o\omega.\beta Q'+p\pi.\beta Q''\\
0&=nv.\gamma R+o\omega.\gamma R'+p\pi.\gamma R'',
\end{aligned}\]
atque formentur hinc istæ æquationes,
\[\begin{aligned}
0&=\alpha P+\beta Q+\gamma R\\
0&=\alpha P'+\beta Q'+\gamma R'\\
0&=\alpha P''+\beta Q''+\gamma R''.
\end{aligned}\]
Hic statim patet, si pro $\alpha,\beta,\gamma$ accipiantur quantitates constantes, tum primam æquationem reliquas binas ultro in se complecti; si enim fuerit $0=\alpha P+\beta Q+\gamma R$, tum simul erit $0=\alpha dP+\beta dQ+\gamma dR$, \& $0=\alpha ddP+\beta ddQ+\gamma ddR$;
\sourcepage{236}\& quia est $P'=P+dP$; $Q'=Q+dQ$, $R'=R+dR$; atque $P''=P+2dP+ddP$; $Q''=Q+2dQ+ddQ$ \& $R''=R+2dR+ddR$; fiet quoque
\[0=\alpha P'+\beta Q'+\gamma R'\quad\&\quad 0=\alpha P''+\beta Q''+\gamma R''.\]
Quocirca ad Problema solvendum formanda est hæc æquatio
\[0=\alpha P+\beta Q+\gamma R;\]
quæ, si loco $\alpha,\beta$, \& $\gamma$ quantitates quæcunque constantes arbitrariæ scribantur, exprimet naturam curvæ quæsitæ. Congruit autem omnino hæc æquatio cum ea quam altera Methodo elicimus, alteraque Methodus per alteram confirmatur. Q. E. I.
\subsection*{Corollarium I.}
\originalsection{16}Omnia ergo hujus quoque generis Problemata resolvi possunt, ope valorum differentialium ex unius applicatæ mutatione oriundorum, quos supra satis ampliter invenire docuimus.
\subsection*{Corollarium II.}
\originalsection{17}Manifestum igitur est, si curva debeat inveniri quæ, inter omnes alias ad eandem abscissam relatas atque in quas binæ expressiones $A$ \& $B$ æqualiter competant, habeat valorem expressionis $C$ maximum minimumve; tum quæstionem redire ad hanc, quæ ad Methodum absolutam pertineat, ut, inter omnes omnino curvas ad eandem abscissam relatas, determinetur ea in qua sit expressio $\alpha A+\beta B+\gamma C$ maximum vel minimum.
\subsection*{Corollarium III.}
\originalsection{18}Simul vero etiam hinc Methodus patet resolvendi Problemata, \sourcepage{237}in quibus, inter omnes curvas in quas plures duabus atque adeo quotcunque proprietates æqualiter conveniant, ea requiritur quæ maximi minimive cujusdam proprietate gaudeat.
\subsection*{Corollarium IV.}
\originalsection{19}Quod si enim, inter omnes curvas in quibus expressiones $A,B,C,D$ æquales obtineant valores, ea debeat investigari in qua sit expressio $E$ maximum vel minimum; tum quæsito satisfiet, si inter omnes omnino curvas ea quæratur in qua sit $\alpha A+\beta B+\gamma C+\delta D+\varepsilon E$ maximum vel minimum; denotantibus litteris $\alpha,\beta,\gamma,\delta,\varepsilon$ quantitates quascunque constantes \& arbitrarias.
\subsection*{Corollarium V.}
\originalsection{20}Quo plures igitur proponantur proprietates, quæ iis curvis ex quibus quæsitam maximi minimive indole præditam indagare oportet, communes esse debeant; eo plures in æquationem pro curva ingredientur quantitates constantes arbitrariæ; atque adeo eo plures curvæ satisfacientes in ea comprehendentur.
\subsection*{Scholion I.}
\originalsection{21}Cur eo plures constantes in Solutionem ingrediantur, quo plures proponantur proprietates communes, ex præcedentibus facile colligi potest. Ponamus enim, inter omnes curvas eadem proprietate $A$ gaudentes, eam investigari oportere, in qua sit $B$ maximum vel minimum: ac primo quidem constabit huic Quæstioni eam curvam esse satisfacturam quæ, inter omnes omnino curvas, habeat $B$ maximum vel minimum; hæc enim, inter omnes quoque illas quæ secum eadem communi proprietate $A$ gaudebunt, habebit $B$ maximum vel minimum. Deinde autem licet innumerabilia istiusmodi curvarum genera concipi, quæ singula eundem valorem expressionis $A$ recipiant; in uno \sourcepage{238}quoque vero genere una erit curva, quæ præ reliquis valorem expressionis $B$ contineat maximum vel minimum. Necesse autem est has curvas satisfacientes omnes in Solutione generali contineri debere. Cum igitur, ob unam proprietatem communem præscriptam numerus curvarum satisfacientium fiat infinitus, multo magis is augebitur, propter eandem rationem, si plures proprietates communes proponantur. Interim tamen si valores, quos habent singulæ proprietates communes in curvis ex quibus quæsitam erui oportet actu definiantur, tum utique solutio unicam Curvam satisfacientem præbebit. Constantes scilicet illæ eo inservient ut valores, quos proprietates communes in curva inventa obtinebunt, pro arbitrio determinentur; sic, per has constantes, in casu duarum proprietatum communium $A$ \& $B$, curva poterit assignari quæ datos expressionum $A$ \& $B$ recipiat valores, atque insuper ita sit comparata ut, inter omnes infinitas alias eosdem illarum expressionum $A$ \& $B$ valores recipientes, habeat valorem alius cujuscunque expressionis $C$ maximum vel minimum. Atque hæc eadem admonitio locum habet, si plures proprietates communes fuerint præscriptæ; ex quo satis perspicuum est, quid hisce constantibus in Solutionem ingredientibus sit faciendum, \& quomodo eas ad usum traduci oporteat: id quod in sequentibus Exemplis clarius declarari poterit.
\subsection*{Exemplum I.}
\originalsection{22}\textit{Inter omnes curvas ad eandem abscissam $AC=a$ relatas, quæ cum inter se ejusdem sint longitudinis, tum etiam æquales areas $DAD$ comprehendant; determinare eam, quæ circa axem $AC$ rotata generet solidum maximæ vel minimæ capacitatis.} (Fig.~14.)

Positis abscissa $AP=x$, applicata $PM=y$ \& $dy=pdx$; binæ proprietates communes propositæ sunt $\int ydx$ \& $\int dx\sqrt{(1+pp)}$; at maximi minimive formula est $\int yydx$. Quærantur jam harum trium formularum valores differentiales. Ac primo quidem erit formulæ $\int ydx$ valor differentialis $=nv.dx$; deinde formulæ
\sourcepage{239}$\int dx\sqrt{(1+pp)}$ valor differentialis est $=-nv.d.\frac p{\sqrt{(1+pp)}}$: \& tertio formulæ $\int yydx$ valor differentialis est $=2nv.ydx$. Ex quibus tribus valoribus differentialibus conficietur pro curva quæsita ista æquatio,
\[0=\alpha dx-\beta d.\frac p{\sqrt{(1+pp)}}+2\gamma ydx,\]
seu
\[ccd.\frac p{\sqrt{(1+pp)}}=bdx+2ydx=\frac{ccdp}{(1+pp)^{3:2}}.\]
Multiplicetur hæc æquatio per $p$ \& integretur; habebitur $ff+by+yy=\frac{-cc}{\sqrt{(1+pp)}}$, ubi tam $cc$ quam $ff$ pro arbitrio sive affirmative sive negative accipere licet. Hinc porro fiet $(ff+by+yy)^2(1+pp)=c^4$; \&
\[p=\frac{\sqrt{(c^4-(ff+by+yy)^2)}}{ff+by+yy}=\frac{dy}{dx};\]
ideoque
\[dx=\frac{(ff+by+yy)dy}{\sqrt{(c^4-(ff+by+yy)^2)}};\]
quæ est æquatio pro curva Elastica. Ingredietur autem, per integrationem unam reliquam, nova quarta constans arbitraria; atque hisce quatuor constantibus effici poterit ut curva per data duo puncta transeat; deinde binis reliquis constantibus obtinebitur, ut posito $x=a$, tam area curvæ quam ejus longitudo datam magnitudinem consequantur. Insuper autem ambiguitate signorum, qua signum radicale afficitur, alterum signum præbebit curvam maximi, alterum minimi proprietate gaudentem. Quoniam autem in æquatione inventa data illa abscissæ magnitudo $a$ non inest; sequitur curvæ inventæ portionem quamvis cuicunque abscissæ respondentem hac quoque gaudere prærogativa ut, inter omnes alias curvas eidem illi abscissæ respondentes, \& per eadem duo puncta transeuntes, quæ simul cum illa curva tum æqualem longitudinem quam æqualem aream complectantur, ut illa, inquam, curva circa abscissam suam rotata generet solidum maximæ minimæve capacitatis. Duo scilicet puncta, per quæ curva quæsita transeat, ideo hic in considerationem sunt ducenda, quia calculus præbuit æquationem differentialem secundi gradus, quæ per se duplicem determinationem requirit. Poterunt vero etiam \sourcepage{240}binæ reliquæ constantes, quæ statim in æquatione inventa inerant, per puncta determinari, hocque pacto determinata solutio hujusmodi emerget, quæ docebit per quatuor data puncta curvam describere, quæ inter omnes alias per eadem quatuor puncta transeuntes, atque cum æque longas tum æquales areas continentes, producat circa axem rotata solidum vel maximum vel minimum. Perpetuo nimirum numerus constantium arbitrariarum, quæ in æquatione inventa cum actu tum potentia insunt, declarabit quot determinationes sint adhibendæ ut curva determinetur; hæcque deinde, inter omnes alias curvas iisdem determinationibus præditas, quæsito satisfaciet.
\subsection*{Exemplum II.}
\originalsection{23}\textit{Inter omnes lineas eidem abscissæ respondentes, quæ primo æquales contineant areas $\int ydx$; atque præterea circum axem rotatæ æqualia generent solida $\int yydx$; determinare eam quæ suum gravitatis centrum vel maxime vel minime habeat elevatum: hoc est in qua sit $\frac{\int yxdx}{\int ydx}$ vel maximum vel minimum.}

Sit abscissæ longitudo præscripta, ad quam Solutionem accommodari oportet, $=a$; atque pro hac abscissa fiat valor formulæ $\int ydx=A$, formulæ $\int yydx=B$, \& formulæ $\int yxdx=C$. Porro sit valor differentialis formulæ $\int ydx=dA=dx$; formulæ $\int yydx=dB=2ydx$, \& formulæ $\int yxdx=dC=xdx$; sumptis nimirum harum formularum valoribus differentialibus secundum regulas supra datas: omittendo tantum particulam $nv$, quippe quæ perpetuo per divisionem tollitur. Cum jam maximi minimive expressio sit, non simplex formula, sed fractio $\frac{\int yxdx}{\int ydx}$; ejus valor differentialis erit
\[\frac{AdC-CdA}{A^2}=\frac{Axdx-Cdx}{A^2};\]
atque, ob proprietatum binarum communium $\int ydx$ \& $\int yydx$ valores differentiales datos, nempe $dA=dx$ \& $dB=2ydx$, resultabit pro curva quæsita sequens æquatio
\sourcepage{241}
\[\alpha dx+2\beta ydx+\frac{\gamma Axdx-\gamma Cdx}{A^2}=0,\]
vel $(\alpha A^2-\gamma C)dx+2\beta A^2ydx+\gamma Axdx=0$; in qua æquatione, cum $\alpha,\beta,\gamma$ sint constantes arbitrariæ, earum transformatione simul constantes determinatæ $A$ \& $C$ ex computo expelli possunt; ita ut Solutio inventa ad omnes abscissas æque fiat accommodata. Pervenietur autem ad hanc æquationem $bdx=mydx+nxdx$, seu, per $dx$ divisione instituta, $b=my+nx$, quæ æquatio est pro linea recta quacunque. Linea recta igitur ad axem verticalem utcunque sita, inter omnes alias lineas cum axe tam eandem aream $\int ydx$ quam idem volumen $\int yydx$ continentes, habebit suæ areæ centrum gravitatis vel maxime vel minime elevatum. Erit autem centrum gravitatis minime elevatum, si linea recta sursum cum axe convergat; maxime autem erit elevatum, si deorsum cum axe convergat; hique sunt ambo casus, quibus vel maxima vel minima centri gravitatis elevatio locum habet. Inter hos casus est medius, quo linea illa recta fit axi parallela: de quo dubium succrescere potest, utrum centrum gravitatis sit vel maxime depressum vel maxime elevatum. Verum iste casus nequidem in Quæstione locum invenit. Nam posita linea recta axi parallela, ita ut sit $y=b$, tum omnino nulla alia exhiberi potest linea quæ pro eadem abscissa, cum æqualem aream $\int ydx$ tum æquale volumen $\int yydx$ contineat: hocque ideo evenit, quod ista linea recta, inter omnes alias lineas eandem aream $\int ydx$ comprehendentes, minimum volumen $\int yydx$ includat.
\subsection*{Exemplum III.}
\originalsection{24}\textit{Inter omnes curvas ejusdem longitudinis $DAD$ puncta data $D,D$ jungentes, determinare eam cujus hæc sit proprietas ut, si inter rectas verticales $DB,DB$ per horizontalem $NN$ spatium $NDADN$ datæ magnitudinis abscindatur, hujus spatii $NDADN$ Centrum gravitatis imum obtineat locum.} (Fig.~21.)

Quæstionis hujus Solutio eximium habet usum in Hydrostatica,
\sourcepage{242}ejusque ope solvetur Problema quo figura lintei $DAD$ vasi $BDDB$ in punctis $DD$ annexi investigatur, quam induit, si vasi data aquæ copia infundatur. Primo enim dum linteum extensionem non admittit, longitudo curvæ $DAD$ erit data, deinde etiam spatium $NDADN$, quo quantitas aquæ infusæ mensuratur, erit data: ac tertio, secundum generales Hydrostaticæ \& gravitationis leges, figuram $DAD$ ita comparatam esse oportet, ut spatii $NDADN$ centrum gravitatis infimum occupet locum. Ad hoc Problema resolvendum, ponatur $DC=CD=a$, \& ducta horizontali quacunque $MPM$, sit $MP=PM=x$, \& $AP=y$, erit arcus $MAM=2\int dx\sqrt{(1+pp)}$, posito $dy=pdx$. Quod si jam longitudo curvæ $DAD$ ponatur $=2b$; æquatio inter $x$ \& $y$ ita debet esse comparata, ut formula integralis $\int dx\sqrt{(1+pp)}$ fiat $=b$, posito $x=a$. Porro area $MAM$ fit $=2\int xdy=2\int xpdx$; quæ fiat $=2ff$, casu quo ponitur $x=a$; ita ut tum sit $\int xpdx=ff$. Hæc vero area non ipsa est data, sed ea cum area $NDDN$ datum spatium producere debet quod sit $=2cc$. Si igitur ponatur $DN=z$, erit $az+ff=cc$, \& $z=\frac{cc-ff}{a}=\frac{cc-\int xpdx}{a}$, posito $x=a$. Denique centrum gravitatis totius spatii $NDADN$ a puncto $A$ distabit intervallo
\[\frac{\int xypdx+az(AC+\frac12z)}{cc},\]
posito post integrationem $x=a$; infra punctum $C$ igitur centrum gravitatis situm erit intervallo
\[\frac{AC(cc-az)-\frac12azz-\int xypdx}{cc},\]
quod debet esse maximum. Cum vero sit $z=\frac{cc-\int xpdx}{a}$; maximum esse debebit hæc forma
\[AC\int xpdx-\frac{c^4}{2a}+\frac{cc\int xpdx}{a}-\frac{(\int xpdx)^2}{2a}-\int xypdx.\]
Problema itaque huc redit ut, inter omnes curvas ejusdem longitudinis datæ abscissæ $DC=a$ respondentes, definiatur ea in qua sit hæc expressio
\[h\int xpdx+\frac{cc}{a}\int xpdx-\frac1{2a}(\int xpdx)^2-\int xypdx\]
\sourcepage{243}maximum, existente $y=h$, posito $x=a$. Jam quia longitudo curvæ est $=\int dx\sqrt{(1+pp)}$, erit ejus valor differentialis $=-d.\frac p{\sqrt{(1+pp)}}$. Deinde formulæ $\int xpdx$ valor differentialis est $=-dx$, \& valor differentialis formulæ $\int xypdx=xpdx-d.xy=-ydx$. Hinc totius expressionis, quæ maximum esse debet, valor differentialis prodit $=-hdx-\frac{cc}{a}dx+\frac{ff}{a}dx+ydx$, quæ, ob $h$ \& $ff$ constantes non determinatas, transit in hanc $kdx+ydx$; ubi $k$ est constans arbitraria. Quocirca prodibit ista æquatio pro curva quæsita
\[kdx+ydx=-ggd.\frac p{\sqrt{(1+pp)}};\]
quæ, per $p$ multiplicata \& integrata, dabit
\[m+2ky+yy=\frac{2gg}{\sqrt{(1+pp)}}:\]
quam curvam constat esse Elasticam, manebitque ea invariata, quemcunque valorem obtineat quantitas $cc$. Quæstioni ergo propositæ ita satisfiet, ut per data puncta $D$ \& $D$ curva Elastica traducatur, cujus axis seu diameter orthogonalis sit recta verticalis $AC$, \& cujus portio $DAD$ datam obtineat longitudinem $2b$; hocque pacto, Solutio omnino erit determinata, unicaque curva satisfaciens resultabit. Quod autem quantitas spatii $NDADN=2cc$, de cujus centro gravitatis quæstio est, prorsus ex computo excesserit, id quidem facile prævidere licuisset; quo pacto, Solutio multo facilior extitisset. Verum data opera hanc conditionem, etsi inutilem, adjecimus, ut modus pateret alia istiusmodi Problemata, ubi talis reductio locum non invenit, resolvendi.
\subsection*{Scholion II.}
\originalsection{25}Sic igitur exposita est universa Methodus maximorum \& minimorum indeterminata, qua linea curva quæri solet maximi minimive proprietate quapiam prædita. Itaque Methodus tota perducta est ad inventionem valorum differentialium, qui ex unius tantum applicatæ incremento oriuntur. Scilicet si Problema \sourcepage{244}postulet, inter omnes omnino curvas ad eandem abscissam relatas, eam in qua expressio quæpiam indefinita maximum minimumve obtineat valorem; tum illius expressionis quærendus est valor differentialis; qui nihilo æqualis positus dabit æquationem pro Curva quæsita. Quod si autem, inter omnes curvas quæ una pluribusve proprietatibus communibus gaudeant, eam definiri oporteat, in qua valor cujuspiam expressionis propositæ fiat maximus vel minimus; tum, tam singularum proprietatum communium, quam maximi minimive, expressionis quæri debent valores differentiales, hique singuli per constantes arbitrarias multiplicari, quorum productorum summa nihilo æqualis posita dabit æquationem pro Curva quæsita. Ad valorem autem differentialem cujusque expressionis indeterminatæ inveniendum, Regulas in superioribus Capitibus sufficientes atque admodum faciles tradidimus. Ejusmodi enim expressio indeterminata, sive proprietatem communem continens, sive maximum minimumve, perpetuo vel est formula integralis simplex, vel functio duarum pluriumve hujusmodi formularum integralium. Quod vero ad formulas integrales simplices attinet; in Cap. IV §. 7 præcepta exposuimus, quorum ope ejusmodi formularum valores differentiales reperiri queant; ubi hanc indagationem ad quinque casus reduximus. Quemadmodum autem secundum hæc eadem præcepta, cujuscunque functionis duarum pluriumve formularum integralium simplicium valor differentialis conveniens definiri queat, id in ejusdem Cap. IV, Propositione 4, indicavimus, modumque differentiationis similem atque satis facilem exposuimus: ita ut in hoc genere nihil superesse videatur, quod insuper sit adjiciendum.
\begin{center}\textsc{Finis.}\end{center}
